\documentclass[a4paper,12pt]{article}

\usepackage[T2A]{fontenc}
\usepackage[utf8]{inputenc}
\usepackage[russian]{babel}

\usepackage{amsmath, amssymb, amsthm, mathtools}
\usepackage{helvet}
\renewcommand{\familydefault}{\sfdefault}
\usepackage{icomma}
\usepackage{geometry}
\geometry{left=2.2cm, right=2.2cm, top=2.5cm, bottom=2.5cm}
\usepackage{microtype}
\usepackage{tikz}
\usetikzlibrary{calc}
\usepackage{pgfplots}
\pgfplotsset{compat=1.18}
\usepackage{tabularx, booktabs, longtable}
\usepackage{enumitem}
\usepackage{hyperref}
\usepackage{parskip}
\usepackage{fancyhdr}
\usepackage{setspace}
\usepackage{xcolor}

\definecolor{accent}{HTML}{008080}
\definecolor{darkaccent}{HTML}{004D4D}
\definecolor{textgray}{HTML}{333333}

\renewcommand{\arraystretch}{1.35}

\pagestyle{fancy}
\fancyhf{}
\fancyhead[L]{\small\color{textgray} Рациональные неравенства}
\fancyhead[R]{\small\color{textgray} Чек-лист}
\fancyfoot[C]{\small\thepage}
\renewcommand{\headrulewidth}{0.4pt}
\renewcommand{\footrulewidth}{0pt}

\fancypagestyle{plain}{%
  \fancyhf{}
  \fancyfoot[C]{\small\thepage}
  \renewcommand{\headrulewidth}{0pt}
}

\newcommand{\sectionheader}[1]{%
  \vspace{0.6cm}%
  \noindent{\large\color{darkaccent}\bfseries #1}%
  \vspace{0.3cm}%
  \par\noindent\rule{\linewidth}{0.4pt}\par\vspace{0.2cm}%
}

\newcommand{\subsectionheader}[1]{%
  \vspace{0.4cm}%
  \noindent{\normalsize\color{accent}\bfseries #1}%
  \par\vspace{0.15cm}%
}

\newenvironment{checklist}{%
  \begin{itemize}[leftmargin=1.5em, itemsep=0.2em, parsep=0.1em]%
}{%
  \end{itemize}%
}

\newcommand{\chk}[1]{\item[\color{accent}\checkmark] #1}

\begin{document}

% ============================================================
% ТИТУЛЬНЫЙ ЛИСТ
% ============================================================
\thispagestyle{empty}
\begin{center}

\vspace*{2cm}

{\LARGE\color{darkaccent}\bfseries Чек-лист}

\vspace{0.6cm}

{\Large\color{accent} Рациональные неравенства}

\vspace{0.3cm}

{\large Метод интервалов, алгебраические приёмы,}\\[0.1cm]
{\large оценка области значений и замена переменных}

\vspace{2.5cm}

{\normalsize Лёвин Артём Александрович}\\[0.2cm]
{\normalsize эксклюзивно для Николь}

\vspace{1cm}

{\normalsize\today}

\vfill

\begin{tikzpicture}
  \draw[color=accent, line width=1.2pt, opacity=0.6]
    (0,0) .. controls (3,1.5) and (6,-0.5) .. (9,0.8)
    .. controls (11,1.5) and (13,0.2) .. (15,0.6);
  \draw[color=darkaccent, line width=0.6pt, opacity=0.35]
    (0,-0.3) .. controls (4,0.8) and (8,-0.8) .. (11,0.3)
    .. controls (13,0.9) and (14,-0.1) .. (15,0.2);
\end{tikzpicture}

\end{center}

\newpage

% ============================================================
% ОСНОВНОЙ ЧЕК-ЛИСТ
% ============================================================
\setcounter{page}{1}

\sectionheader{1. Базовый алгоритм метода интервалов}

\begin{checklist}
  \chk \textbf{Стандартный вид:} неравенство должно быть приведено к виду $\frac{P(x)}{Q(x)} \lor 0$, где $\lor$ --- это один из знаков $>, <, \ge, \le$. Справа всегда должен стоять ноль.
  \chk \textbf{Нули числителя:} решаем уравнение $P(x) = 0$.
  \chk \textbf{Нули знаменателя (ОДЗ):} решаем уравнение $Q(x) = 0$. Знаменатель не может обращаться в ноль.
  \chk \textbf{Отметка на числовой прямой:}
  \begin{itemize}
    \item Корни знаменателя отмечаем \textbf{выколотой} (пустой) точкой.
    \item Корни числителя отмечаем \textbf{закрашенной} точкой, если неравенство нестрогое ($\ge, \le$), и \textbf{выколотой}, если строгое ($>, <$).
  \end{itemize}
  \chk \textbf{Расстановка знаков:} выбираем удобное значение $x$ (например, $x=0$) из любого интервала, подставляем в исходное выражение и определяем знак. Далее знаки чередуются при переходе через корни нечётной кратности.
  \chk \textbf{Запись ответа:} выбираем интервалы с нужным знаком. Границы интервалов записываются в соответствии с типом точки (квадратная скобка для закрашенной, круглая для выколотой).
\end{checklist}

\subsectionheader{Пример}
Решить неравенство: $\dfrac{x+1}{x-1} \ge 0$.

\textit{Шаг 1.} Корни: числитель $x = -1$, знаменатель $x = 1$.

\textit{Шаг 2.} На числовой прямой $x = -1$ --- закрашенная (неравенство нестрогое), $x = 1$ --- выколотая (ОДЗ).

\begin{center}
\begin{tikzpicture}[scale=0.8]
  \draw[->, thick, textgray] (-4,0) -- (4,0);
  \fill[accent] (-1.5,0) circle (3pt) node[below] {$-1$};
  \draw[fill=white, draw=accent, thick] (1.5,0) circle (3pt) node[below] {$1$};
  \node[above, textgray] at (-2.5,0) {$+$};
  \node[above, textgray] at (0,0) {$-$};
  \node[above, textgray] at (2.5,0) {$+$};
  \draw[thick, accent] (-4,0) -- (-1.5,0);
  \draw[thick, accent] (1.5,0) -- (4,0);
\end{tikzpicture}
\end{center}

\textbf{Ответ:} $(-\infty;\,-1] \cup (1;\,+\infty)$.

\sectionheader{2. Приведение к стандартному виду}

\begin{checklist}
  \chk \textbf{Запрет на умножение:} категорически запрещено умножать обе части неравенства на знаменатель (или любое выражение с переменной). Это приводит к потере ОДЗ и ошибкам в определении знака.
  \chk \textbf{Перенос и общий знаменатель:} все слагаемые переносятся в одну часть (обычно влево), после чего выражение приводится к общему знаменателю.
  \chk \textbf{Упрощение числителя:} после приведения к общему знаменателю раскрываем скобки в числителе и приводим подобные. Знаменатель оставляем в факторизованном виде (не раскрываем скобки, чтобы не потерять корни).
\end{checklist}

\sectionheader{3. Кратные корни и изолированные точки}

\begin{checklist}
  \chk \textbf{Чётная кратность:} если корень встречается чётное число раз (например, $(x-a)^2$ или $(x-b)^4$), то при переходе через него на числовой прямой \textbf{знак не меняется}.
  \chk \textbf{Изолированные точки:} если корень чётной кратности принадлежит числителю, а неравенство нестрогое, и при этом сам корень не входит в ОДЗ (или не попадает в нужный интервал), он включается в ответ как отдельное изолированное число (через объединение $\cup \{a\}$).
\end{checklist}

\subsectionheader{Пример}
Решить неравенство: $\dfrac{(x-7)^2}{x(x-3)} \le 0$.

\textit{Анализ:} Корни числителя: $x=7$ (кратность 2). Корни знаменателя: $x=0, x=3$.
Знак не меняется при переходе через $7$. Точка $7$ --- закрашенная, $0$ и $3$ --- выколотые.

\begin{center}
\begin{tikzpicture}[scale=0.6]
  \draw[->, thick, textgray] (-1,0) -- (10,0);
  \draw[fill=white, draw=accent, thick] (1,0) circle (3pt) node[below] {$0$};
  \draw[fill=white, draw=accent, thick] (4,0) circle (3pt) node[below] {$3$};
  \fill[accent] (8,0) circle (3pt) node[below] {$7$};
  \node[above, textgray] at (0,0) {$+$};
  \node[above, textgray] at (2.5,0) {$-$};
  \node[above, textgray] at (6,0) {$+$};
  \node[above, textgray] at (9,0) {$+$};
  \draw[thick, accent] (1,0) -- (4,0);
\end{tikzpicture}
\end{center}

\textbf{Ответ:} $(0;\,3) \cup \{7\}$.

\sectionheader{4. Алгебраические приёмы: группировка и замена}

\begin{checklist}
  \chk \textbf{Метод группировки:} применяется, когда в числителе или знаменателе 4 слагаемых. Группируем первое со вторым и третье с четвёртым (или первое с четвёртым, второе с третьим), выносим общие множители, затем выносим общую скобку.
  \chk \textbf{Метод замены:} применяется при наличии повторяющихся блоков (например, $x^4$ и $x^2$). Вводим замену $t = x^2$.
  \chk \textbf{Ограничение замены:} так как $x^2 \ge 0$ для любых действительных $x$, то и $t \ge 0$. Отрицательные значения $t$ отбрасываем.
\end{checklist}

\sectionheader{5. Метод оценки (область значений)}

\begin{checklist}
  \chk \textbf{Сумма чётных степеней:} если выражение имеет вид $x^{2n} + x^{2k} + \dots + C$, где все коэффициенты положительны и $C > 0$, то такое выражение \textbf{строго больше нуля} при любых $x$.
  \chk \textbf{Упрощение:} положительный множитель в числителе или знаменателе не влияет на знак дроби. Его можно мысленно «отсечь», но если он был в знаменателе, ОДЗ на него всё равно не накладывается (так как он никогда не равен нулю).
\end{checklist}

\sectionheader{6. Типичные ошибки и ловушки}

\begin{checklist}
  \chk \textbf{Сокращение дробей до метода интервалов:} при сокращении одинаковых скобок в числителе и знаменателе (например, $\frac{x-3}{(x-1)(x-3)}$) легко забыть, что $x \neq 3$. Корень знаменателя всегда должен быть выколот на числовой прямой.
  \chk \textbf{Потеря изолированных точек:} забывают включить в ответ корни чётной кратности, если они не попадают в основной интервал.
  \chk \textbf{Игнорирование ОДЗ при замене:} забывают проверить условие $t \ge 0$ при замене $t = x^2$ или $t = |x|$.
  \chk \textbf{Умножение на знаменатель:} приводит к потере знака и получению неверного ответа.
\end{checklist}

\newpage

% ============================================================
% ТРЕНИРОВОЧНЫЙ БЛОК
% ============================================================

\sectionheader{Тренировочный блок}

\noindent\textit{Решите рациональные неравенства. Упражнения расположены по возрастанию сложности.}

\vspace{0.4cm}

\noindent\textbf{Уровень: средний}

\vspace{0.2cm}

\noindent\textbf{1.}\quad
$\dfrac{x-3}{x+2} \ge 0$

\vspace{0.3cm}

\noindent\textbf{2.}\quad
$\dfrac{3x+6}{x-4} < 0$

\vspace{0.3cm}

\noindent\textbf{3.}\quad
$\dfrac{x^2-9}{x^2-4x+3} \le 0$

\vspace{0.5cm}

\noindent\textbf{Уровень: выше среднего}

\vspace{0.2cm}

\noindent\textbf{4.}\quad
$\dfrac{x^2-5x+6}{x^2+x-12} > 0$

\vspace{0.3cm}

\noindent\textbf{5.}\quad
$\dfrac{3x-1}{x+2} - 3 \le 0$

\vspace{0.3cm}

\noindent\textbf{6.}\quad
$\dfrac{(x-4)^2}{x^2-9} \ge 0$

\vspace{0.3cm}

\noindent\textbf{7.}\quad
$\dfrac{x^3-x^2-4x+4}{x+3} \le 0$

\vspace{0.3cm}

\noindent\textbf{8.}\quad
$\dfrac{x^4+3x^2+2}{x^2-6x+8} < 0$

\vspace{0.3cm}

\noindent\textbf{9.}\quad
$\dfrac{x^4-5x^2+4}{x^2+x+1} \ge 0$

\vspace{0.5cm}

\noindent\textbf{Уровень: сложный}

\vspace{0.2cm}

\noindent\textbf{10.}\quad
$\dfrac{x^4-8x^2+16}{x^2-5x+6} \le 0$

\newpage

% ============================================================
% ОТВЕТЫ
% ============================================================

\thispagestyle{fancy}

\sectionheader{Ответы к тренировочному блоку}

\vspace{0.3cm}

\begin{center}
\renewcommand{\arraystretch}{1.5}
\begin{tabularx}{\linewidth}{@{}c X@{}}
\toprule
\textbf{\textnumero} & \textbf{Ответ} \\
\midrule
1 & $(-\infty;\,-2) \cup [3;\,+\infty)$ \\
2 & $(-2;\,4)$ \\
3 & $[-3;\,1)$ \\
4 & $(-\infty;\,-4) \cup (2;\,3) \cup (3;\,+\infty)$ \\
5 & $(-2;\,+\infty)$ \\
6 & $(-\infty;\,-3) \cup (3;\,+\infty) \cup \{4\}$ \\
7 & $(-3;\,-2] \cup [1;\,2]$ \\
8 & $(2;\,4)$ \\
9 & $(-\infty;\,-2] \cup [-1;\,1] \cup [2;\,+\infty)$ \\
10 & $\{-2\} \cup (2;\,3)$ \\
\bottomrule
\end{tabularx}
\end{center}

\vspace{1cm}

\noindent\textit{Примечание к задаче 3.} После факторизации получаем $\frac{(x-3)(x+3)}{(x-1)(x-3)} \le 0$. Точка $x=3$ исключается из ОДЗ, поэтому она не входит в ответ, несмотря на то, что формально обнуляет числитель.

\noindent\textit{Примечание к задаче 6.} Корень $x=4$ имеет чётную кратность в числителе. Знак дроби при переходе через него не меняется, но так как неравенство нестрогое ($\ge 0$), сама точка включается в ответ как изолированная.

\noindent\textit{Примечание к задаче 8.} Числитель $x^4+3x^2+2 > 0$ при любых $x$ (метод оценки). Неравенство сводится к анализу знака знаменателя.

\noindent\textit{Примечание к задаче 10.} Числитель сворачивается в $(x^2-4)^2 = (x-2)^2(x+2)^2$. Знаменатель равен $(x-2)(x-3)$. Точка $x=2$ является корнем знаменателя и обязательно выкалывается. Точка $x=-2$ обнуляет числитель и входит в ОДЗ, поэтому включается как изолированная. На интервале $(2;\,3)$ знаменатель отрицателен, а числитель положителен, что удовлетворяет условию $\le 0$.

\end{document}